\documentclass[10pt,a4paper]{amsart}
\usepackage[margin=19mm]{geometry}
\usepackage{amsmath,amssymb,mathtools}
\usepackage[scr=rsfs,cal=euler]{mathalfa}
\usepackage{xfrac}
\usepackage[unicode,hidelinks,bookmarks=false]{hyperref}
\newcommand{\ie}{i.e.\ }
\newcommand{\cf}{cf.\ }
\newcommand{\resp}{respectively\ }
\newcommand{\Subsection}{\subsection}
\providecommand{\nobreakdash}{\nobreak\hbox{-}}
\setlength{\emergencystretch}{2em}
\multlinegap0pt
\usepackage{amssymb}
\usepackage{xcolor}
\usepackage[shortlabels]{enumitem}
\newtheorem{proposition}{Proposition}
\newtheorem{lemma}{Lemma}
\title{\textbf{Scattering of the defocusing Calogero--Moser derivative nonlinear Schrödinger equation}}
\author{Xi Chen}
\date{Source: arXiv:2511.06432v6 (CC BY 4.0)}
\begin{document}
\maketitle

\noindent\emph{Source and licence.} \textbf{Scattering of the defocusing Calogero--Moser derivative nonlinear Schrödinger equation}. Version arXiv:2511.06432v6 (CC BY 4.0), distributed by arXiv under the Creative Commons Attribution 4.0 International licence, http://creativecommons.org/licenses/by/4.0/, as recorded in the arXiv licence metadata for this exact version and shown on the abstract page https://arxiv.org/abs/2511.06432v6. Full licence notice: \texttt{licenses/arxiv-2511.06432-CC-BY-4.0.txt}.\par\smallskip

\noindent\textit{Editorial scope.} Counted unit: the article's lemma a.7 (source0.tex lines 6815--6902). The article's complete original proof of it is delivered with it (source0.tex lines 6904--7538, one \texttt{proof} environment of 635 lines). No mathematics is written, repaired or re-derived: the statement and the proof are the article's own lines, in the article's own notation, and no retained line is substituted or re-flowed.  Retained uncounted as the source-local context the proof needs: lines 4001--4045; lines 6750--6769.  Nothing was omitted from any retained span, so the package carries no omission record.  No retained line contains a citation, so the wrapper carries no bibliography and no bibliography entry.  No retained line contains a figure environment, an image-inclusion command or drawing code, so no image is delivered and the package declares no asset dependency.  Editorial formatting only, in addition: an AMS \texttt{amsart} preamble that restates the article's own declarations and macros used by the retained lines, namely the environments \texttt{proposition}, \texttt{lemma} on separate counters, together with the standard packages those lines need.  The retained environments print in this document as Proposition 1, Lemma 1 in order of appearance, whereas the article prints the same items with the numbering of its own full document.  The complete native source of the article as distributed in the arXiv e-print for this exact version is bundled unchanged as \texttt{sources/188729/source0.tex}.
\begin{proposition}[The action of \(X^*\) under the distorted Fourier transform]
\label{prop:Xstar_under_Fd}
Let \(u\in\mathcal S_+(\mathbb R)\), and set
\[
\mathfrak u:=\widetilde u:=\mathcal F_{d,u}u.
\]
Let
\[
\widetilde f\in H^1(0,\infty),
\qquad
f:=\mathcal F_{d,u}^{-1}\widetilde f.
\]
Then
\[
f\in\operatorname{Dom}(X^*),
\qquad
I_+(f)=J_+(\widetilde f):=\widetilde f(0_{+}),
\]
and
\[
\mathcal F_{d,u}(X^*f)
=
i\partial_\lambda\widetilde f
+
B(\mathfrak u)\widetilde f,
\]
where
\[
B(\mathfrak u)g
:=
-\frac1{4\pi}
\left(
i\mathfrak u\,H(\overline{\mathfrak u}g)
-
|\mathfrak u|^2g
\right).
\]
Here \(H\) denotes the Hilbert transform on the positive half-line
\[
H h(\mu)
:=
\frac1\pi\operatorname{p.v.}
\int_0^\infty\frac{h(\lambda)}{\mu-\lambda}\,d\lambda,\quad \mu >0.
\]
\end{proposition}

\begin{proposition}
\label{prop 2}
Let $z\in\mathbb C_{+}$. We have
\begin{equation}\label{eq:L1-L2-bound 4}
\forall f \in L^2(0,\infty), \qquad \|(i\partial_\lambda-z)^{-1}f\|_{L^2(0,\infty)}
\leq \frac{1}{\Im z}\,\|f\|_{L^2(0,\infty)}.
\end{equation}
\begin{equation}\label{eq:L1-L2-bound 1}
\forall f \in L^1(0,\infty), \qquad \|(i\partial_\lambda-z)^{-1}f\|_{L^2(0,\infty)}
\leq \frac{1}{\sqrt{2\,\Im z}}\,\|f\|_{L^1(0,\infty)}.
\end{equation}
\begin{equation}\label{eq:L1-L2-bound 2}
\forall f \in L^2(0,\infty), \qquad \|(i\partial_\lambda-z)^{-1}f\|_{L^\infty(0,\infty)}
\leq \frac{1}{\sqrt{2\,\Im z}}\,\|f\|_{L^2(0,\infty)}.
\end{equation}
\begin{equation}\label{eq:L1-L2-bound 3}
\forall f \in L^1(0,\infty), \qquad\|(i\partial_\lambda-z)^{-1}f\|_{L^\infty(0,\infty)}
\leq \|f\|_{L^1(0,\infty)}.
\end{equation}
\end{proposition}

\begin{lemma}[The rough boundary functional]
\label{coro a.7}
Let
\[
\mathfrak u\in L^2(0,\infty).
\]
For \(w\in L^2(0,\infty)\) and \(r\in L^\infty(0,\infty)\), define
\[
B(w)r
:=
-\frac1{4\pi}
\left(
iw\,H(\overline w r)-|w|^2r
\right),
\]
where \(H\) denotes the Hilbert transform on the positive half-line,
\[
Hh(\mu)
=
\frac1\pi\operatorname{p.v.}
\int_0^\infty
\frac{h(\lambda)}{\mu-\lambda}\,d\lambda .
\]
For smooth \(w\in\mathcal S_*\), set
\[
A(w,t,z)
:=
\left(
i\partial_\lambda
+
e^{-it\lambda^2}B(w)e^{it\lambda^2}
-z
\right)^{-1},
\qquad z\in\mathbb C_+ .
\]
Then there exists a family of bounded linear functionals
\[
\mathcal J_{\mathfrak u}(t,z):L^1(0,\infty)\to\mathbb C,
\qquad t\in\mathbb R,\quad z\in\mathbb C_+,
\]
with the following properties.

First, for every \(\gamma>0\), there exists \(C_{\mathfrak u,\gamma}>0\) such that
\begin{equation}
\label{A7-boundary-L1-bound}
|\mathcal J_{\mathfrak u}(t,z)[f]|
\le
C_{\mathfrak u,\gamma}\|f\|_{L^1(0,\infty)}
\end{equation}
for all \(t\in\mathbb R\), all \(z\in\mathbb C_+\) with
\[
\Im z\ge\gamma,
\]
and all \(f\in L^1(0,\infty)\). The constant is independent of \(t\) and
\(\Re z\).

Second, if
\[
\mathfrak u_n\in \mathcal S_*,
\qquad
\mathfrak u_n\to\mathfrak u
\quad\text{in }L^2(0,\infty),
\]
and if
\[
f_n\in L^1(0,\infty)\cap L^2(0,\infty),
\qquad
f_n\to f
\quad\text{in }L^1(0,\infty),
\]
then, for every fixed \(t\in\mathbb R\) and \(z\in\mathbb C_+\),
\begin{equation}
\label{A7-boundary-approx-conv}
J_+\bigl[A(\mathfrak u_n,t,z)f_n\bigr]
\longrightarrow
\mathcal J_{\mathfrak u}(t,z)[f],
\end{equation}
where \(J_+\) is defined in Proposition \(\ref{prop:Xstar_under_Fd}\).

We shall use the notation
\begin{equation}
\label{A7-boundary-notation}
J_+\bigl[A(\mathfrak u,t,z)f\bigr]
:=
\mathcal J_{\mathfrak u}(t,z)[f],
\qquad f\in L^1(0,\infty).
\end{equation}
\end{lemma}

\begin{proof}
We write
\[
R_z:=(i\partial_\lambda-z)^{-1},
\qquad
\mathcal C_H:=\frac1{4\pi}(\operatorname{Id}-iH).
\]
Thus
\[
B(w)=M_w\mathcal C_HM_{\overline w}
\]
whenever both sides are defined. Indeed,
\[
M_w\mathcal C_HM_{\overline w}g
=
\frac1{4\pi}
\left(
|w|^2g-iwH(\overline wg)
\right)
=
-\frac1{4\pi}
\left(
iwH(\overline wg)-|w|^2g
\right).
\]
For \(w\in L^2(0,\infty)\), set
\[
M_{w,t}:=M_{e^{-it\lambda^2}w},
\qquad
M_{\overline w,t}:=M_{\overline w\,e^{it\lambda^2}}.
\]
For smooth \(w\), we have
\[
e^{-it\lambda^2}B(w)e^{it\lambda^2}
=
M_{w,t}\mathcal C_HM_{\overline w,t}.
\]
We also introduce
\begin{equation}
\label{A7-Gamma-def}
\Gamma_w(t,z)
:=
\mathcal C_HM_{\overline w,t}R_zM_{w,t}.
\end{equation}

We first record that \(B(w)r\in L^1(0,\infty)\) whenever
\(w\in L^2(0,\infty)\) and \(r\in L^\infty(0,\infty)\). Since \(H\) is bounded
on \(L^2(0,\infty)\),
\[
\begin{aligned}
\|B(w)r\|_{L^1}
&\le
C\|wH(\overline wr)\|_{L^1}
+
C\||w|^2r\|_{L^1}       \\
&\le
C\|w\|_{L^2}\|H(\overline wr)\|_{L^2}
+
C\|w\|_{L^2}^2\|r\|_{L^\infty}  \\
&\le
C\|w\|_{L^2}^2\|r\|_{L^\infty}.
\end{aligned}
\]

\medskip
\noindent
\textbf{Step 1: Basic \(L^2\) and Hilbert--Schmidt estimates.}
For \(z=x+iy\in\mathbb C_+\), the free resolvent is given by
\[
(R_z f)(\lambda) : = (i\partial_\lambda-z)^{-1} f
=
i\int_\lambda^\infty e^{iz(s-\lambda)}f(s)\,ds .
\]
By Proposition \ref{prop 2}, we have
\begin{equation}
\label{A7-Rz-basic-bounds}
\|R_z\|_{L^1\to L^\infty}\le1,
\qquad
\|R_z\|_{L^1\to L^2}\le (2y)^{-1/2},
\end{equation}
and
\begin{equation}
\label{A7-Rz-L2-bounds}
\|R_z\|_{L^2\to L^\infty}\le (2y)^{-1/2},
\qquad
\|R_z\|_{L^2\to L^2}\le y^{-1}.
\end{equation}
For \(a,b\in L^2(0,\infty)\), the operator
\[
M_{\overline a,t}R_zM_{b,t}
\]
is Hilbert--Schmidt on \(L^2(0,\infty)\), and
\begin{equation}
\label{A7-HS-mixed}
\begin{aligned}
\|M_{\overline a,t}R_zM_{b,t}\|_{\mathrm{HS}}^2
&=
\int_0^\infty\int_\lambda^\infty
|a(\lambda)|^2 e^{-2y(s-\lambda)}|b(s)|^2\,ds\,d\lambda  \\
&\le
\|a\|_{L^2}^2\|b\|_{L^2}^2 .
\end{aligned}
\end{equation}
In particular,
\[
\|M_{\overline a,t}R_zM_{b,t}\|_{\mathcal L(L^2,L^2)}
\le
\|a\|_{L^2}\|b\|_{L^2}.
\]
For \(w\in L^2(0,\infty)\), define
\[
\Theta_w(y)^2
:=
\int_0^\infty\int_\lambda^\infty
|w(\lambda)|^2 e^{-2y(s-\lambda)}|w(s)|^2\,ds\,d\lambda .
\]
Then
\[
\Theta_w(y)\to0
\qquad\text{as }y\to+\infty.
\]
Indeed, the integrand is dominated by
\[
|w(\lambda)|^2|w(s)|^2\mathbf 1_{\{s>\lambda\}},
\]
which is integrable on \((0,\infty)^2\), and it converges pointwisely to \(0\) as
\(y\to+\infty\).

Since \(\mathcal C_H\) is bounded on \(L^2(0,\infty)\), \eqref{A7-HS-mixed}
implies
\begin{equation}
\label{A7-Gamma-HS-bound}
\|\Gamma_w(t,z)\|_{\mathcal L(L^2,L^2)}
\le
C\Theta_w(\Im z).
\end{equation}
Moreover, if \(w_n\to w\) in \(L^2(0,\infty)\), then for each fixed
\(z\in\mathbb C_+\),
\begin{equation}
\label{A7-Gamma-L2-continuity}
\Gamma_{w_n}(t,z)\to\Gamma_w(t,z)
\quad\text{in }\mathcal L(L^2,L^2),
\end{equation}
uniformly in \(t\) and \(\Re z\). Indeed,
\[
\Gamma_{w_n}-\Gamma_w
=
\mathcal C_HM_{\overline{w_n-w},t}R_zM_{w_n,t}
+
\mathcal C_HM_{\overline w,t}R_zM_{w_n-w,t},
\]
and \eqref{A7-HS-mixed} gives
\[
\|\Gamma_{w_n}-\Gamma_w\|_{\mathcal L(L^2,L^2)}
\le
C\|w_n-w\|_{L^2}
\bigl(\|w_n\|_{L^2}+\|w\|_{L^2}\bigr).
\]

\medskip
\noindent
\textbf{Step 2: Large-\(\Im z\) construction for \(L^1\)-inputs.}
Choose \(b>0\) so large that
\begin{equation}
\label{A7-large-im-smallness}
\|\Gamma_{\mathfrak u}(t,z)\|_{\mathcal L(L^2,L^2)}
\le
\frac14
\qquad
\text{whenever }\Im z\ge b.
\end{equation}
The choice of \(b\) is independent of \(t\) and \(\Re z\). If
\[
\mathfrak u_n\to\mathfrak u
\quad\text{in }L^2(0,\infty),
\]
then, after possibly increasing \(n\), the same estimate with \(1/2\) in place
of \(1/4\) holds for \(\mathfrak u_n\), uniformly for \(\Im z\ge b\).

For \(\Im z\ge b\) and \(f\in L^1(0,\infty)\), define
\begin{equation}
\label{A7-A1-large-def}
A^{(1)}(\mathfrak u,t,z)f
:=
R_zf
-
R_zM_{\mathfrak u,t}
\bigl(I+\Gamma_{\mathfrak u}(t,z)\bigr)^{-1}
\mathcal C_HM_{\overline{\mathfrak u},t}R_zf .
\end{equation}
Then
\[
A^{(1)}(\mathfrak u,t,z):L^1(0,\infty)\to L^2(0,\infty)
\]
is bounded. Indeed, by \eqref{A7-Rz-basic-bounds},
\[
\|R_zf\|_{L^2}\le (2\Im z)^{-1/2}\|f\|_{L^1},
\]
and
\[
\|M_{\overline{\mathfrak u},t}R_zf\|_{L^2}
\le
\|\mathfrak u\|_{L^2}\|R_zf\|_{L^\infty}
\le
\|\mathfrak u\|_{L^2}\|f\|_{L^1}.
\]
Thus
\[
h:=
\bigl(I+\Gamma_{\mathfrak u}(t,z)\bigr)^{-1}
\mathcal C_HM_{\overline{\mathfrak u},t}R_zf
\]
belongs to \(L^2\), with
\[
\|h\|_{L^2}\le C_{\mathfrak u}\|f\|_{L^1}.
\]
Since
\[
\|M_{\mathfrak u,t}h\|_{L^1}
\le
\|\mathfrak u\|_{L^2}\|h\|_{L^2},
\]
another use of \eqref{A7-Rz-basic-bounds} gives
\begin{equation}
\label{A7-A1-large-bound}
\|A^{(1)}(\mathfrak u,t,z)f\|_{L^2}
\le
C_{\mathfrak u,b}\|f\|_{L^1},
\qquad \Im z\ge b.
\end{equation}

For \(\Im z\ge b\), define the large-\(\Im z\) boundary functional by
\begin{equation}
\label{A7-J1-large-def}
\mathcal J^{(1)}_{\mathfrak u}(t,z)[f]
:=
J_+[R_zf]
-
J_+\left[
R_zM_{\mathfrak u,t}
\bigl(I+\Gamma_{\mathfrak u}(t,z)\bigr)^{-1}
\mathcal C_HM_{\overline{\mathfrak u},t}R_zf
\right].
\end{equation}
The first term is
\[
J_+[R_zf]
=
i\int_0^\infty e^{iz\lambda}f(\lambda)\,d\lambda,
\]
and hence
\begin{equation}
\label{A7-free-boundary-L1}
|J_+[R_zf]|
\le
\|f\|_{L^1}.
\end{equation}
For the second term, if \(h\in L^2(0,\infty)\), then
\begin{equation}
\label{A7-JRMu-bound}
|J_+[R_zM_{\mathfrak u,t}h]|
\le
\|M_{\mathfrak u,t}h\|_{L^1}
\le
\|\mathfrak u\|_{L^2}\|h\|_{L^2}.
\end{equation}
Therefore
\begin{equation}
\label{A7-J1-large-bound}
|\mathcal J^{(1)}_{\mathfrak u}(t,z)[f]|
\le
C_{\mathfrak u,b}\|f\|_{L^1},
\qquad \Im z\ge b.
\end{equation}
The constants above are independent of \(t\) and \(\Re z\).

For smooth \(w\in\mathcal S_*\), the Woodbury formula gives
\[
A(w,t,z)f
=
R_zf
-
R_zM_{w,t}
\bigl(I+\Gamma_w(t,z)\bigr)^{-1}
\mathcal C_HM_{\overline w,t}R_zf
\]
for \(f\in L^1(0,\infty)\cap L^2(0,\infty)\). Hence for \(f\in L^1(0,\infty)\cap L^2(0,\infty)\) and $\Im z \ge b$, 
\[
J_+[A(w,t,z)f]
=
\mathcal J^{(1)}_w(t,z)[f].
\]

\medskip
\noindent
\textbf{Step 3: The auxiliary \(L^2\)-input boundary functionals.}
For \(\Im z\ge b\) and \(g\in L^2(0,\infty)\), define
\begin{equation}
\label{A7-A2-large-def}
A^{(2)}(\mathfrak u,t,z)g
:=
R_zg
-
R_zM_{\mathfrak u,t}
\bigl(I+\Gamma_{\mathfrak u}(t,z)\bigr)^{-1}
\mathcal C_HM_{\overline{\mathfrak u},t}R_zg ,
\end{equation}
and
\begin{equation}
\label{A7-J2-large-def}
\mathcal J^{(2)}_{\mathfrak u}(t,z)[g]
:=
J_+[R_zg]
-
J_+\left[
R_zM_{\mathfrak u,t}
\bigl(I+\Gamma_{\mathfrak u}(t,z)\bigr)^{-1}
\mathcal C_HM_{\overline{\mathfrak u},t}R_zg
\right].
\end{equation}
By \eqref{A7-Rz-L2-bounds},
\[
\|M_{\overline{\mathfrak u},t}R_zg\|_{L^2}
\le
(2\Im z)^{-1/2}\|\mathfrak u\|_{L^2}\|g\|_{L^2},
\]
and hence
\[
A^{(2)}(\mathfrak u,t,z):L^2(0,\infty)\to L^2(0,\infty)
\]
is bounded for \(\Im z\ge b\). Moreover,
\[
|J_+[R_zg]|
\le
\left(\int_0^\infty e^{-2(\Im z)\lambda}\,d\lambda\right)^{1/2}
\|g\|_{L^2}
=
(2\Im z)^{-1/2}\|g\|_{L^2},
\]
and the second term in \(\mathcal J^{(2)}_{\mathfrak u}\) is controlled by
\eqref{A7-JRMu-bound}. Thus
\[
\mathcal J^{(2)}_{\mathfrak u}(t,z)\in (L^2(0,\infty))^*
\]
for \(\Im z\ge b\).

For smooth \(w\in\mathcal S_*\), the smooth theory gives that
\[
A(w,t,z)
=
\left(
i\partial_\lambda+
e^{-it\lambda^2}B(w)e^{it\lambda^2}
-z
\right)^{-1}
\]
is the resolvent of a maximally dissipative operator. Hence
\begin{equation}
\label{A7-smooth-resolvent-bound}
\|A(w,t,z)\|_{\mathcal L(L^2,L^2)}
\le
\frac1{\Im z},
\qquad z\in\mathbb C_+,
\end{equation}
and the resolvent identity holds:
\begin{equation}
\label{A7-smooth-resolvent-identity}
A(w,t,z)
=
A(w,t,z_0)
\bigl[I-(z-z_0)A(w,t,z_0)\bigr]^{-1}
\end{equation}
In particular, if
\[
|z-z_0|<\frac12\Im z_0,
\]
then the inverse in \eqref{A7-smooth-resolvent-identity} is uniformly bounded.

In the large-\(\Im z\) region, the explicit formulae
\eqref{A7-A2-large-def} and \eqref{A7-J2-large-def}, together with
\eqref{A7-Gamma-L2-continuity}, imply that if
\[
w_n\to\mathfrak u
\qquad\text{in }L^2(0,\infty),
\]
then
\begin{equation}
\label{A7-A2-large-conv}
A(w_n,t,z)\to A^{(2)}(\mathfrak u,t,z)
\quad\text{in }\mathcal L(L^2,L^2),
\qquad \Im z\ge b,
\end{equation}
and
\begin{equation}
\label{A7-J2-large-conv}
J_+[A(w_n,t,z)\,\cdot]
\to
\mathcal J^{(2)}_{\mathfrak u}(t,z)
\quad\text{in }(L^2(0,\infty))^*,
\qquad \Im z\ge b.
\end{equation}

We now continue the \(L^2\)-input objects to all of \(\mathbb C_+\). Suppose that
at some \(z_0\in\mathbb C_+\) we have already constructed
\(A^{(2)}(\mathfrak u,t,z_0)\) and
\(\mathcal J^{(2)}_{\mathfrak u}(t,z_0)\), and that for every smooth
approximating sequence \(w_n\to\mathfrak u\) in \(L^2\),
\[
A(w_n,t,z_0)\to A^{(2)}(\mathfrak u,t,z_0)
\quad\text{in }\mathcal L(L^2,L^2),
\]
and
\[
J_+[A(w_n,t,z_0)\,\cdot]
\to
\mathcal J^{(2)}_{\mathfrak u}(t,z_0)
\quad\text{in }(L^2(0,\infty))^*.
\]
Since \eqref{A7-smooth-resolvent-bound} passes to the limit, we have
\[
\|A^{(2)}(\mathfrak u,t,z_0)\|_{\mathcal L(L^2,L^2)}
\le
\frac1{\Im z_0}.
\]
Thus, if
\[
|z-z_0|<\frac12\Im z_0,
\]
we may define
\begin{equation}
\label{A7-A2-continuation}
A^{(2)}(\mathfrak u,t,z)
:=
A^{(2)}(\mathfrak u,t,z_0)
\bigl[
I-(z-z_0)A^{(2)}(\mathfrak u,t,z_0)
\bigr]^{-1},
\end{equation}
and
\begin{equation}
\label{A7-J2-continuation}
\mathcal J^{(2)}_{\mathfrak u}(t,z)[g]
:=
\mathcal J^{(2)}_{\mathfrak u}(t,z_0)
\left[
\bigl(
I-(z-z_0)A^{(2)}(\mathfrak u,t,z_0)
\bigr)^{-1}g
\right].
\end{equation}
The smooth resolvent identity \eqref{A7-smooth-resolvent-identity} and the
convergence at \(z_0\) imply the corresponding convergence at \(z\). Therefore
the definition can be propagated along any finite chain of such discs.

Starting from the large-\(\Im z\) region and using a finite chain of discs, we
obtain
\[
\mathcal J^{(2)}_{\mathfrak u}(t,z)\in (L^2(0,\infty))^*
\]
for every \(z\in\mathbb C_+\). The definition is independent of the approximating
sequence and of the chain, because it is obtained as the norm limit of the smooth
functionals
\[
J_+[A(w_n,t,z)\,\cdot],
\]
and this limit is unique.

Moreover, for every \(\gamma>0\), the same finite-step continuation argument yields
\begin{equation}
\label{A7-J2-uniform-bound}
|\mathcal J^{(2)}_{\mathfrak u}(t,z)[g]|
\le
C^{(2)}_{\mathfrak u,\gamma}\|g\|_{L^2}
\end{equation}
for all \(t\in\mathbb R\), all \(z\in\mathbb C_+\) with
\(\Im z\ge\gamma\), and all \(g\in L^2(0,\infty)\). The constant is independent
of \(t\) and \(\Re z\).

\medskip
\noindent
\textbf{Step 4: Extension of the \(L^1\)-input boundary functional to all \(z\).}
Fix \(\gamma>0\). Choose \(Y\ge b\), and let
\[
z=x+iy\in\mathbb C_+,
\qquad y\ge\gamma.
\]
If \(y\ge Y\), we use the large-\(\Im z\) definition
\eqref{A7-J1-large-def}. If \(\gamma\le y<Y\), set
\[
z_Y:=x+iY.
\]
For \(f\in L^1(0,\infty)\), define
\begin{equation}
\label{A7-J1-lower-def}
\mathcal J_{\mathfrak u}(t,z)[f]
:=
\mathcal J^{(1)}_{\mathfrak u}(t,z_Y)[f]
+
(z-z_Y)
\mathcal J^{(2)}_{\mathfrak u}(t,z)
\left[
A^{(1)}(\mathfrak u,t,z_Y)f
\right].
\end{equation}
This is meaningful because
\[
A^{(1)}(\mathfrak u,t,z_Y)f\in L^2(0,\infty)
\]
by \eqref{A7-A1-large-bound}, and
\(\mathcal J^{(2)}_{\mathfrak u}(t,z)\) is bounded on \(L^2\) by
\eqref{A7-J2-uniform-bound}. Combining \eqref{A7-J1-large-bound},
\eqref{A7-A1-large-bound}, and \eqref{A7-J2-uniform-bound}, we obtain
\[
|\mathcal J_{\mathfrak u}(t,z)[f]|
\le
C_{\mathfrak u,\gamma}\|f\|_{L^1},
\]
uniformly for \(t\in\mathbb R\), \(\Re z\in\mathbb R\), and \(\Im z\ge\gamma\).
This proves \eqref{A7-boundary-L1-bound}.

The definition is independent of the auxiliary height \(Y\). Indeed, if
\(f\in L^1\cap L^2\) and \(w\in\mathcal S_*\), then the smooth resolvent identity
gives
\[
A(w,t,z)
=
A(w,t,z_Y)
+
(z-z_Y)A(w,t,z)A(w,t,z_Y),
\]
and therefore \eqref{A7-J1-lower-def} agrees with the usual smooth boundary
value
\[
J_+[A(w,t,z)f].
\]
Passing to smooth approximations \(w_n\to\mathfrak u\) in \(L^2\), we obtain the
same consistency for rough \(\mathfrak u\) and \(f\in L^1\cap L^2\). Since
\(L^1\cap L^2\) is dense in \(L^1\), and both possible definitions are bounded
\(L^1\)-functionals with the same type of bound, the definition is independent
of \(Y\) for all \(f\in L^1\).

\medskip
\noindent
\textbf{Step 5: Stability under approximation.}
Let
\[
\mathfrak u_n\in\mathcal S_*,
\qquad
\mathfrak u_n\to\mathfrak u
\quad\text{in }L^2(0,\infty),
\]
and let
\[
f_n\in L^1(0,\infty)\cap L^2(0,\infty),
\qquad
f_n\to f
\quad\text{in }L^1(0,\infty).
\]
Fix \(t\in\mathbb R\) and \(z=x+iy\in\mathbb C_+\). Choose \(Y\ge b\) so large
that the large-\(\Im z\) construction is valid for \(\mathfrak u\) and for all
sufficiently large \(n\), and set
\[
z_Y:=x+iY.
\]

In the large-\(\Im z\) region, the explicit formulae
\eqref{A7-A1-large-def} and \eqref{A7-J1-large-def}, together with
\[
\mathfrak u_n\to\mathfrak u
\quad\text{in }L^2,
\qquad
f_n\to f
\quad\text{in }L^1,
\]
give
\begin{equation}
\label{A7-A1-approx-conv}
A^{(1)}(\mathfrak u_n,t,z_Y)f_n
\longrightarrow
A^{(1)}(\mathfrak u,t,z_Y)f
\qquad\text{in }L^2(0,\infty),
\end{equation}
and
\begin{equation}
\label{A7-J1-approx-large}
\mathcal J^{(1)}_{\mathfrak u_n}(t,z_Y)[f_n]
\longrightarrow
\mathcal J^{(1)}_{\mathfrak u}(t,z_Y)[f].
\end{equation}
For instance,
\[
\|M_{\overline{\mathfrak u_n},t}R_{z_Y}f_n
-
M_{\overline{\mathfrak u},t}R_{z_Y}f\|_{L^2}
\le
\|\mathfrak u_n-\mathfrak u\|_{L^2}\|R_{z_Y}f\|_{L^\infty}
+
\|\mathfrak u_n\|_{L^2}\|R_{z_Y}(f_n-f)\|_{L^\infty},
\]
and the right-hand side tends to zero. The remaining terms are handled in the
same way, using \eqref{A7-Gamma-L2-continuity} and the convergence of the
inverses.

By Step 3,
\[
\mathcal J^{(2)}_{\mathfrak u_n}(t,z)
\to
\mathcal J^{(2)}_{\mathfrak u}(t,z)
\]
in the operator norm of functionals on \(L^2(0,\infty)\), and these functionals
are uniformly bounded. Combining this with \eqref{A7-A1-approx-conv}, we get
\[
\mathcal J^{(2)}_{\mathfrak u_n}(t,z)
\left[
A^{(1)}(\mathfrak u_n,t,z_Y)f_n
\right]
\to
\mathcal J^{(2)}_{\mathfrak u}(t,z)
\left[
A^{(1)}(\mathfrak u,t,z_Y)f
\right].
\]
Using \eqref{A7-J1-lower-def} for \(\mathfrak u_n\) and for \(\mathfrak u\), and
then using \eqref{A7-J1-approx-large}, we obtain
\[
\mathcal J_{\mathfrak u_n}(t,z)[f_n]
\to
\mathcal J_{\mathfrak u}(t,z)[f].
\]
For smooth \(\mathfrak u_n\), the left-hand side is exactly
\[
J_+\bigl[A(\mathfrak u_n,t,z)f_n\bigr].
\]
Thus \eqref{A7-boundary-approx-conv} follows, and the proof is complete.
\end{proof}

\end{document}
